\chapter{Выходы машины: как мозг управляет мышцами}
\chaptermark{Выходы машины}
\label{sec:muscles}

\section{Быстрый выход}

Всё, что машина делает в мире быстро, она делает мышцами. Слово, взгляд, шаг, улыбка --- это сокращения мышц. Второй, медленный выход --- химия тела --- разобран в главе~\ref{sec:chemoutput}. На рис.~\ref{fig:machine} выход был стрелкой «мышцы»; теперь посмотрим, что за этой стрелкой.

Последнее звено цепи --- \nt{ACET}-нейроны, те самые, у которых в таблице~\ref{tab:types} записано «выход на мышцу». Каждое мышечное волокно получает вход ровно от одного такого нейрона, а один нейрон управляет группой волокон --- от нескольких сотен до пары тысяч~\cite{Feinstein1955mu}. В мышцах, которым нужна тонкая подстройка, единицы мельче, в больших мышцах ног --- крупнее. Нейрон вместе со своими волокнами называют \emph{двигательной единицей}: это наименьший кусок мышцы, который мозг может включить отдельно.

Синапс нейрона на мышце устроен совсем не так, как синапсы коры из главы~\ref{sec:code}. Там большая часть импульсов терялась. Здесь каждый импульс выбрасывает в несколько раз больше ацетилхолина, чем нужно, чтобы волокно сработало~\cite{WoodSlater2001}. Это синапс с \emph{запасом надёжности}: один импульс нейрона --- ровно одно сокращение волокон. Внутри машины можно позволить себе ненадёжные связи и исправлять ошибки избыточностью; на выходе ошибка стоит движения, и надёжность куплена прямо в железе.

\section{Сила: частота и число единиц}

Один импульс даёт одиночное сокращение --- подёргивание: сила нарастает за несколько десятков миллисекунд и спадает. Хорошее приближение для него~\cite{Fuglevand1993}
\begin{equation}
h(t) \;=\; F_1\,\frac{t}{T_c}\,e^{\,1-t/T_c},
\label{eq:twitch}
\end{equation}
где $T_c$ --- время нарастания, порядка $50\ms$, а $F_1$ --- сила одного подёргивания. Если импульсы идут чаще, подёргивания складываются:
\begin{empheq}[box=\keybox]{equation}
F(t) \;=\; \sum_k h\bigl(t-t_k\bigr) .
\label{eq:forcesum}
\end{empheq}
Это тот же фильтр нижних частот, что превращал импульсы в концентрацию в главе~\ref{sec:serocomputer}, только теперь на выходе не концентрация, а сила. Мышца --- цифро-аналоговый преобразователь из импульсов в силу (рис.~\ref{fig:tetanus}). В нашей модели при пяти импульсах в секунду подёргивания почти не перекрываются и сила скачет от $0.2F_1$ до $1.1F_1$; при пятнадцати она уже держится между $1.8F_1$ и $2.2F_1$; при сорока становится почти ровной, около $5.4F_1$. Настоящая мышца при частых импульсах насыщается, так что линейная сумма~\eqref{eq:forcesum} верна лишь до нескольких десятков импульсов в секунду.

\begin{figure}[t]
\centering
\begin{tikzpicture}[>=Stealth,x=20cm,y=0.55cm]
\draw[->,gray!70!black] (0,0) -- (0.66,0) node[right,font=\small,black] {$t$, с};
\draw[->,gray!70!black] (0,0) -- (0,6.4) node[left,font=\small,black] {$F/F_1$};
\foreach \t in {0,0.2,0.4,0.6}{\draw[gray!60!black] (\t,-0.1) -- (\t,0.1); \node[below,font=\scriptsize] at (\t,0) {\t};}
\foreach \f in {1,3,5}{\draw[gray!60!black] (-0.004,\f) -- (0.004,\f); \node[left,font=\scriptsize] at (0,\f) {\f};}
\draw[thick,blue!60!black] plot file {figures/muscle_twitch_5.dat};
\draw[thick,green!45!black] plot file {figures/muscle_twitch_15.dat};
\draw[thick,red!70!black] plot file {figures/muscle_twitch_40.dat};
\node[font=\scriptsize,blue!60!black,anchor=west] at (0.605,0.9) {5 имп/с};
\node[font=\scriptsize,green!45!black,anchor=west] at (0.605,2.1) {15 имп/с};
\node[font=\scriptsize,red!70!black,anchor=west] at (0.605,5.4) {40 имп/с};
\end{tikzpicture}
\caption{Мышца как цифро-аналоговый преобразователь: сила~\eqref{eq:forcesum} при регулярных импульсах одной двигательной единицы, $T_c=50\ms$. Редкие импульсы дают отдельные подёргивания, частые сливаются в ровную силу --- так же, как частые импульсы сливаются в ровную концентрацию на рис.~\ref{fig:lowpass}.}
\label{fig:tetanus}
\end{figure}

У мозга две ручки силы: частота импульсов каждой единицы и число включённых единиц. И вторая ручка устроена очень изобретательно. Единицы в мышце разные: самые слабые в сотню раз слабее самых сильных. Когда сила нужна всё большая, единицы включаются \emph{по порядку размера}, от маленьких к большим~\cite{Henneman1957}. Никто этот порядок не вычисляет: маленькие нейроны просто легче возбуждаются от того же общего входа, так что порядок включения задан самим железом.

Что это даёт? Пусть каждая следующая по порядку единица сильнее предыдущей в $q$ раз, где $q$ немного больше единицы. Сила всех включённых единиц --- сумма геометрической прогрессии, и почти вся она приходится на последние единицы. Тогда очередная включённая единица добавляет
\begin{empheq}[box=\keybox]{equation}
\frac{\Delta F}{F} \;\approx\; q-1 \;=\; \text{const} ,
\label{eq:sizeprinciple}
\end{empheq}
то есть шаг силы пропорционален самой силе. При слабом усилии мозг дозирует его мелкими порциями, при сильном --- крупными. Это \emph{число с плавающей точкой}: порядок задаётся тем, сколько единиц включено, а мантисса --- частотой их импульсов. Та же логарифмическая шкала, что у датчиков главы~\ref{sec:sensors}: вход и выход машины меряют мир в относительных долях.

У плавающей точки есть обратная сторона. Разброс силы тоже растёт пропорционально ей самой: сильное движение неизбежно менее точно, чем слабое. По одной влиятельной гипотезе, именно этот шум, зависящий от сигнала, объясняет, почему наши движения плавные: плавная команда даёт наименьший разброс в конечной точке~\cite{HarrisWolpert1998}.

\section{Тянуть, но не толкать: два провода на сустав}

Мышца умеет только тянуть. Толкать она не может, как не может \nt{GLUT}-нейрон выдать отрицательный сигнал. Решение знакомо по главе~\ref{sec:glutcomputer}: \emph{два провода}. Каждый сустав обслуживает пара мышц, тянущих в противоположные стороны (рис.~\ref{fig:antagonists}). Если их силы $F_1$ и $F_2$, а плечо $\ell$, то вращающий момент и жёсткость сустава равны
\begin{empheq}[box=\keybox]{equation}
M \;=\; \ell\,(F_1-F_2), \qquad K \;\approx\; \kappa_m\,(F_1+F_2),
\label{eq:antagonists}
\end{empheq}
где $\kappa_m$ --- коэффициент, связывающий силу мышцы с её жёсткостью: напряжённая мышца пружинит сильнее расслабленной. Разность сил задаёт момент, сумма --- жёсткость~\cite{Hogan1984}. Двумя проводами мозг управляет двумя независимыми величинами: куда повернуть сустав и насколько твёрдо его держать. Когда надо удержать чашку, не расплескав, мы напрягаем обе мышцы сразу: момент тот же, жёсткость выше, и случайный толчок меньше сбивает руку.

\begin{figure}[t]
\centering
\begin{tikzpicture}[>=Stealth,
  nrn/.style={circle,draw,very thick,minimum size=0.95cm,inner sep=0pt,font=\scriptsize},
  inh/.style={-{Bar[width=7pt]},thick,red!70!black}]
% the joint: a bar on a pivot, two springs pulling down on either side
\fill[gray!30] (4.6,-0.25) -- (5.4,-0.25) -- (5,0.15) -- cycle;
\draw[line width=3pt,gray!50!black] (2.4,0.2) -- (7.6,0.2);
\fill[black] (5,0.2) circle (0.08);
\draw[decorate,decoration={coil,segment length=5pt,amplitude=4pt},very thick,blue!60!black] (3.0,0.2) -- (3.0,-1.6);
\draw[decorate,decoration={coil,segment length=5pt,amplitude=4pt},very thick,orange!85!black] (7.0,0.2) -- (7.0,-1.6);
\draw[very thick] (2.5,-1.6) -- (3.5,-1.6); \draw[very thick] (6.5,-1.6) -- (7.5,-1.6);
\node[font=\small,blue!60!black] at (2.35,-0.75) {$F_1$};
\node[font=\small,orange!85!black] at (7.65,-0.75) {$F_2$};
\draw[->,thick] (5.9,0.75) arc (60:120:1.8) node[above,font=\scriptsize,pos=0.5] {$M=\ell(F_1-F_2)$};
% motor neurons and reflex circuit
\node[nrn,fill=green!12] (M1) at (1.6,-3.0) {\nt{ACET}};
\node[nrn,fill=green!12] (M2) at (8.4,-3.0) {\nt{ACET}};
\node[nrn,fill=red!10] (G) at (5.0,-3.0) {\nt{GLYC}};
\node[nrn,fill=blue!6] (S) at (1.6,-5.0) {\nt{GLUT}};
\draw[->,very thick] (M1) -- (3.0,-1.7);
\draw[->,very thick] (M2) -- (7.0,-1.7);
\draw[->,thick] (S) -- (M1);
\draw[->,thick] (S) -- (G);
\draw[inh] (G) -- (M2);
\draw[->,thick,densely dashed,gray!60!black] (3.15,-1.2) to[bend left=25] node[pos=0.35,right,font=\scriptsize,black] {растяжение} (S.north east);
\draw[->,thick,blue!60!black] (-0.3,-3.0) node[left,font=\scriptsize,black,align=right] {команда\\мозга} -- (M1);
\draw[->,thick,blue!60!black] (10.3,-3.0) node[right,font=\scriptsize,black,align=left] {команда\\мозга} -- (M2);
\end{tikzpicture}
\caption{Две мышцы на один сустав и самая быстрая петля обратной связи. \nt{ACET}-нейроны получают команды мозга и тянут сустав в разные стороны; разность сил поворачивает его, сумма делает жёстче. Датчик растяжения (\nt{GLUT}) левой мышцы возбуждает её собственный \nt{ACET}-нейрон и через \nt{GLYC}-посредника тормозит нейрон мышцы-противника --- запрет из главы~\ref{sec:gaba}.}
\label{fig:antagonists}
\end{figure}

\section{Обратная связь с задержкой}

Мышцы не работают вслепую. В каждой сидят датчики растяжения, о которых шла речь в таблице~\ref{tab:sensors}, и самая короткая петля замыкается прямо в спинном мозге (рис.~\ref{fig:antagonists}). Мышцу неожиданно растянули --- датчик растяжения возбуждает её собственный \nt{ACET}-нейрон, мышца сокращается и возвращает сустав. Одновременно через тормозящий нейрон-посредник выключается мышца-противник, чтобы не мешала. Этим посредником работает \nt{GLYC} --- тот самый тип из таблицы~\ref{tab:types}, которому не досталось своей главы: его место именно здесь, в быстрых петлях спинного мозга.

У каждой петли есть задержка $d$: у спинномозговой около $25$--$30\ms$ для руки, у петли через кору --- в полтора--три раза больше, у петли через глаз --- $100$--$200\ms$. А задержанная обратная связь, как мы видели в утверждении~\ref{prop:ei}, раскачивает систему. Для простейшего регулятора, исправляющего отклонение $x$ со скоростью $G$ по сведениям $d$ секунд давности, $\dd x/\dd t=-G\,x(t-d)$, условие устойчивости таково~\cite{Hayes1950}:
\begin{empheq}[box=\keybox]{equation}
G\,d \;<\; \frac{\pi}{2} ,
\label{eq:delaystable}
\end{empheq}
а на границе система колеблется с частотой $1/(4d)$. Мы проверили это на модели: при $d=30\ms$ отклонение затухает при $G=40$ в секунду, а при $G=55$ в секунду, чуть выше границы $52$, раскачивается.

Отсюда два следствия. Первое: чем длиннее петля, тем медленнее она обязана исправлять. Через глаз с задержкой в полторы сотни миллисекунд руку нельзя поправлять быстрее, чем примерно за десятую долю секунды, иначе она задрожит. Второе: граница устойчивости спинномозговой петли, $1/(4\cdot30\ms)\approx8$ колебаний в секунду, близка к частоте мелкой дрожи, которая есть у каждого здорового человека, --- восемь-двенадцать колебаний в секунду. Петля растяжения --- один из вероятных источников этой дрожи~\cite{McAuleyMarsden2000}.

Как тогда мозг двигается быстро и точно, если обратная связь запаздывает? По распространённой гипотезе --- \emph{предсказывает}: держит внутреннюю модель тела и вычисляет, каким будет результат команды, не дожидаясь датчиков~\cite{WolpertGhahramani2000}. Датчики нужны, чтобы исправлять модель, а не каждое движение. Инженер назвал бы это управлением с упреждением и наблюдателем состояния.

\section{Генераторы ритма движений}

Ходьба, дыхание, жевание --- ритмические движения. Нужны ли для них часы? Нет. В спинном мозге и стволе есть небольшие сети, которые сами порождают ритм, даже если на них не приходит ничего ритмического~\cite{MarderBucher2001}. Простейшая такая сеть собрана из деталей, которые у нас уже есть: две группы \nt{GLUT}-нейронов, которые тормозят друг друга через \nt{GABA}- или \nt{GLYC}-посредников, как на рис.~\ref{fig:wta}, и медленно устают, как память главы~\ref{sec:glutcomputer}:
\begin{empheq}[box=\keybox]{equation}
\tau\,\frac{\dd r_i}{\dd t} = -r_i + \bigl[I - \beta\,r_j - a_i\bigr]_+ ,
\qquad
\tau_a\,\frac{\dd a_i}{\dd t} = -a_i + b\,r_i ,
\label{eq:halfcentre}
\end{empheq}
где $a_i$ --- усталость группы $i$, $j$ --- другая группа. Взаимное торможение с $\beta>1$ выбирает победителя (утверждение~\ref{prop:wta}). Победитель работает и устаёт; когда усталость съедает его вход, проигравший, уже отдохнувший, вырывается вперёд и сам начинает уставать. Группы работают по очереди, и каждая управляет своей мышцей пары (рис.~\ref{fig:halfcentre}).

\begin{figure}[t]
\centering
\begin{tikzpicture}[>=Stealth,x=3.2cm,y=2.4cm]
\draw[->,gray!70!black] (0,0) -- (4.2,0) node[right,font=\small,black] {$t$, с};
\draw[->,gray!70!black] (0,0) -- (0,1.2) node[left,font=\small,black] {$r$};
\foreach \t in {0,1,2,3,4}{\draw[gray!60!black] (\t,-0.03) -- (\t,0.03); \node[below,font=\scriptsize] at (\t,0) {\t};}
\draw[thick,blue!60!black] plot file {figures/halfcentre1.dat};
\draw[thick,orange!85!black] plot file {figures/halfcentre2.dat};
\node[font=\scriptsize,blue!60!black,anchor=west] at (1.6,1.28) {группа 1: «вперёд»};
\node[font=\scriptsize,orange!85!black,anchor=west] at (2.7,1.28) {группа 2: «назад»};
\end{tikzpicture}
\caption{Генератор ритма по модели~\eqref{eq:halfcentre}: $I=1$, $\beta=2$, $b=2$, $\tau=20\ms$, $\tau_a=0.4\,\text{с}$. Две группы со взаимным торможением и усталостью работают по очереди с периодом $0.84\,\text{с}$, хотя на вход подаётся постоянный сигнал.}
\label{fig:halfcentre}
\end{figure}

В нашей модели при постоянном входе группы сменяют друг друга с периодом $0.84\,\text{с}$, примерно вдвое больше времени усталости $\tau_a$. Постоянная команда сверху --- «иди» --- превращается в ритм на месте. Это ещё один местный ритм, как ритм петли \nt{GLUT}--\nt{GABA} в главе~\ref{sec:gaba}: он возникает, только когда сеть работает, и задаёт такт лишь тем мышцам, которыми управляет.

\begin{keyblock}
\begin{proposition}[Выход машины]
\label{prop:muscles}
Мозг управляет мышцами как иерархия регуляторов. Наверху выбирается действие (глава~\ref{sec:dofa}); ниже местные генераторы превращают постоянные команды в ритм, а быстрые петли спинного мозга держат суставы. Сила задаётся в плавающей точке: порядок --- числом единиц, включаемых по размеру, мантисса --- частотой импульсов. Каждый сустав управляется парой тянущих мышц, разность сил которых задаёт момент, а сумма --- жёсткость. Эти устройства измерены; что мозг опирается на внутреннюю модель тела, --- хорошо обоснованная гипотеза.
\end{proposition}
\end{keyblock}

\section{Итог}

\begin{table}[t]
\centering
\footnotesize
\setlength{\tabcolsep}{4pt}
\renewcommand{\arraystretch}{1.15}
\begin{tabular}{>{\raggedright}p{3.0cm}>{\raggedright}p{4.6cm}>{\raggedright\arraybackslash}p{5.2cm}}
\toprule
Что делает выход & Как & Что известно \\
\midrule
превращает импульсы в силу & сумма подёргиваний~\eqref{eq:forcesum}, фильтр нижних частот & измерено; линейная сумма --- упрощение \\
дозирует силу & включение единиц по размеру, плавающая точка~\eqref{eq:sizeprinciple} & порядок включения измерен \\
толкает, умея только тянуть & пара мышц: момент --- разность, жёсткость --- сумма~\eqref{eq:antagonists} & измерено \\
держит сустав & петля растяжения, запрет противника через \nt{GLYC} & измерено \\
не дрожит & $Gd<\pi/2$~\eqref{eq:delaystable} & следует из теории регулирования; вклад в дрожь --- вероятная причина \\
двигается быстро & предсказание по внутренней модели & хорошо обоснованная гипотеза \\
ходит и дышит ритмично & генератор ритма~\eqref{eq:halfcentre} & генераторы измерены; простейшая схема --- модель \\
\bottomrule
\end{tabular}
\caption{Как машина управляет мышцами.}
\label{tab:muscles}
\end{table}

Выход машины устроен теми же приёмами, что и всё остальное (таблица~\ref{tab:muscles}): два провода вместо знака, фильтр нижних частот вместо ЦАП, логарифмическая шкала вместо линейной, взаимное торможение и усталость вместо тактового генератора. Вход (глава~\ref{sec:sensors}) и выход меряют мир в относительных долях, а между ними работает машина, которой посвящена эта книга.

\section{Задачи}

\begin{exercise}[\lvlA]
\label{ex:13:1}
\emph{Плавающая точка в цифрах.} В мышце $N=100$ двигательных единиц с силами $f_n=f_1q^{\,n-1}$; самая сильная в $100$ раз сильнее самой слабой. (а)~Найдите $q$, относительный шаг~\eqref{eq:sizeprinciple} и динамический диапазон в децибелах. (б)~Каков шаг при силе $10f_1$ у «линейного ЦАП» из $100$ одинаковых единиц той же суммарной силы?
\end{exercise}

\begin{exercise}[\lvlA]
\label{ex:13:2}
\emph{Цена запаса надёжности.} На каждый импульс синапс на мышечном волокне выбрасывает пуассоновское число порций со средним $m$, а волокно срабатывает при $n_{\text{пор}}=20$ порциях и больше; запас надёжности $S=m/n_{\text{пор}}$. Сколько сбоев в сутки даёт единица, работающая с частотой $20$ импульсов в секунду, при $S=2$ и при $S=4$?
\end{exercise}

\begin{exercise}[\lvlB]
\label{ex:13:3}
\emph{Мышца как фильтр.} Подёргивание~\eqref{eq:twitch} с $T_c=50\ms$ --- импульсная характеристика линейной системы. (а)~Найдите её передаточную функцию $H(s)$ и частоту среза. (б)~При какой частоте регулярных импульсов размах пульсаций силы составляет $10\%$ от средней силы?
\end{exercise}

\begin{exercise}[\lvlB]
\label{ex:13:4}
\emph{Быстрее задержки не исправить.} Регулятор $\dd x/\dd t=-G\,x(t-d)$. (а)~Покажите, что быстрее всего отклонение затухает без колебаний при $Gd=1/e$, со скоростью $1/d$. (б)~Найдите это $G$ и постоянную времени затухания для спинномозговой петли, $d=30\ms$, и для петли через глаз, $d=150\ms$.
\end{exercise}

\begin{exercise}[\lvlB]
\label{ex:13:5}
\emph{Период генератора ритма.} При $\tau\ll\tau_a$ период $T$ модели~\eqref{eq:halfcentre} задаётся уравнением $(\beta-1)(1-E^{2+b})/(\beta-E)=b\,(1-E^{1+b})/(1+b)$, где $E=e^{-T/(2\tau_a)}$. (а)~Найдите $T$ при $\beta=b=2$, $\tau_a=0.4$~с. (б)~Покажите, что $T$ не зависит от входа $I$, а ритм возможен только при $b>\beta-1$.
\end{exercise}

\begin{exercise}[\lvlB]
\label{ex:13:6}
\emph{Моделирование генератора.} Импортируйте функцию \texttt{halfcentre} из \texttt{models/\allowbreak muscle\_models.py} (сам файл не запускайте) и измерьте период, отбросив первые два цикла, при $b=0.8$, $1.05$, $1.2$, $1.5$, $2$, $4$ и при $I=0.5$, $1$, $2$. Может ли команда «иди быстрее» менять темп через $I$?
\end{exercise}

\begin{exercise}[\lvlB]
\label{ex:13:7}
\emph{Жёсткость против рефлекса.} Сустав охвачен петлёй растяжения: $\dd\theta/\dd t=-a\theta(t)-G\theta(t-d)$, $a=K/B$, $d=30\ms$. Сведите её к задаче~\ref{ex:13:4} и найдите наименьшее $a$, при котором отклонение затухает без колебаний с постоянной времени $15\ms$, и нужное $G$. Как менять $G$ при усилении совместного напряжения мышц?
\end{exercise}

\begin{exercise}[\lvlC]
\label{ex:13:8}
\emph{Ручка темпа.} По задачам~\ref{ex:13:5} и~\ref{ex:13:6} вход $I$ генератора~\eqref{eq:halfcentre} меняет только размах, но не темп. Предложите минимальное изменение модели из деталей книги, при котором ниже некоторого $I_{\min}$ ритма нет, а выше период убывает с ростом $I$, хотя бы вдвое при утроении $I$. Найдите $I_{\min}$ при $\tau\ll\tau_a$ и проверьте моделированием.
\end{exercise}
